% ===========================================================================
%  Jàng Afrig — Lab 4 : Layouts + Responsive mobile
%  Adapter Jàng Afrig au téléphone avec media queries
%  Licence 1 DUT1-Informatique — ESP/UCAD — 2025–2026
%  Dr. El Hadji Bassirou TOURÉ
%  Compilé avec : xelatex (2–3 passes)
% ===========================================================================
\documentclass[11pt,a4paper]{article}

\usepackage{fontspec}
\setmainfont{Latin Modern Roman}
\setsansfont{Latin Modern Sans}
\setmonofont{Latin Modern Mono}
\usepackage{polyglossia}
\setdefaultlanguage{french}

\usepackage[margin=2.5cm]{geometry}
\usepackage{amsmath,amssymb}
\usepackage{graphicx}
\graphicspath{{figures/}}
\usepackage{xcolor}
\usepackage{booktabs}
\usepackage{hyperref}
\usepackage{tcolorbox}
\tcbuselibrary{theorems,breakable,skins}
\usepackage{enumitem}
\usepackage{fancyhdr}
\usepackage{array}
\usepackage{pifont}
\usepackage{longtable}
\usepackage{multirow}

% --- TikZ ---
\usepackage{tikz}
\usetikzlibrary{shapes.geometric,arrows.meta,positioning,calc,backgrounds,fit,shadows.blur}

\newcommand{\cmark}{\ding{51}}
\newcommand{\xmark}{\ding{55}}

\hypersetup{
    colorlinks=true,
    linkcolor=blue!60!black,
    urlcolor=blue!60!black,
    citecolor=blue!60!black
}

\newtcolorbox{objectifbox}[1][]{colback=blue!5,colframe=blue!50,fonttitle=\bfseries,title=#1,breakable}
\newtcolorbox{infobox}[1][]{colback=green!5,colframe=green!50!black,fonttitle=\bfseries,title=#1,breakable}
\newtcolorbox{attentionbox}[1][]{colback=orange!5,colframe=orange!70,fonttitle=\bfseries,title=#1,breakable}
\newtcolorbox{retenirbox}[1][]{colback=yellow!10,colframe=yellow!50!black,fonttitle=\bfseries,title=#1,breakable}
\newtcolorbox{prereqbox}[1][]{colback=cyan!5,colframe=cyan!50!black,fonttitle=\bfseries,title=#1,breakable}
\newtcolorbox{projetbox}[1][]{colback=teal!5,colframe=teal!50!black,fonttitle=\bfseries,title=#1,breakable}

% --- Couleurs (palette panafricaine) ---
\definecolor{panvert}{HTML}{0E8A4F}
\definecolor{panjaune}{HTML}{F5C518}
\definecolor{panrouge}{HTML}{C8332E}
\definecolor{pannoir}{HTML}{1A1A1A}
\definecolor{panblanc}{HTML}{F8F8F8}

% --- Couleurs schéma breakpoints ---
\definecolor{bpmobile}{RGB}{200,80,60}
\definecolor{bptablet}{RGB}{220,180,90}
\definecolor{bpdesktop}{RGB}{45,160,110}

% --- Code source ---
\usepackage{listings}
\lstdefinelanguage{HTML5}{
  morekeywords={<!DOCTYPE,html,head,meta,title,body,h1,h2,h3,h4,p,ul,ol,li,table,tr,th,td,a,img,strong,em,br,div,span,header,nav,main,section,article,footer,aside,form,input,select,option,textarea,label,button,fieldset,legend,style,link,class,id,href,src,alt,name,content,viewport,figure,figcaption},
  sensitive=true,
  morecomment=[s]{<!--}{-->},
  morestring=[b]",
}
\lstdefinelanguage{CSS}{
  morekeywords={color,background-color,background,font-family,font-size,font-weight,text-align,padding,margin,border,border-radius,width,height,max-width,min-width,display,flex,flex-direction,justify-content,align-items,flex-wrap,gap,box-sizing,transition,transform,position,top,bottom,left,right,opacity,cursor,line-height,letter-spacing,box-shadow,list-style,text-decoration,text-transform,border-collapse,vertical-align,object-fit,overflow,grid,grid-template-columns,grid-gap,media,screen,min-width,max-width,only},
  sensitive=true,
  morecomment=[s]{/*}{*/},
  morestring=[b]",
}
\lstset{
  basicstyle=\ttfamily\small,
  backgroundcolor=\color{gray!8},
  frame=single,
  framerule=0.5pt,
  rulecolor=\color{gray!40},
  breaklines=true,
  breakatwhitespace=true,
  showstringspaces=false,
  tabsize=2,
  xleftmargin=4pt,
  xrightmargin=4pt,
  aboveskip=8pt,
  belowskip=8pt,
  literate={é}{{\'e}}1 {è}{{\`e}}1 {ê}{{\^e}}1 {à}{{\`a}}1 {ù}{{\`u}}1
           {ô}{{\^o}}1 {î}{{\^i}}1 {ç}{{\c{c}}}1 {É}{{\'E}}1
           {→}{{$\rightarrow$}}1
}

\pagestyle{fancy}
\fancyhf{}
\fancyhead[L]{\small Jàng Afrig --- Lab 4 : Responsive mobile}
\fancyhead[R]{\small UCAD --- L1 DUT1-Info}
\fancyfoot[C]{\thepage}
\fancyfoot[L]{\small Dr. El Hadji Bassirou TOURÉ}
\fancyfoot[R]{\small 2025--2026}
\renewcommand{\headrulewidth}{0.4pt}
\renewcommand{\footrulewidth}{0.4pt}

\title{%
\vspace{-1cm}
\textbf{\LARGE Fondamentaux du Web}\\
\textbf{\LARGE HTML, CSS, JavaScript}\\[0.8cm]
\textbf{\huge Lab 4 --- Layouts \& Responsive}\\[0.3cm]
\textbf{\Large Faire en sorte que Jàng Afrig marche sur téléphone}\\[0.5cm]
\large Licence 1 DUT1-Informatique --- 2025--2026}
\author{Dr. El Hadji Bassirou TOURÉ\\
Département Génie Informatique\\
Université Cheikh Anta Diop de Dakar}
\date{}

\begin{document}

\maketitle
\thispagestyle{empty}

\vspace{0.5cm}

% ===========================================================================
% PRÉREQUIS + OBJECTIFS
% ===========================================================================
\begin{prereqbox}[Informations pratiques]
\begin{center}
\begin{tabular}{ll}
\toprule
\textbf{Type} & Lab guidé \\
\textbf{Séance} & S6 --- deuxième Lab CSS \\
\textbf{Prérequis} & Lab 3 terminé (5 pages stylisées avec tag \texttt{v0.3}) \\
\textbf{Validation} & App lisible sur téléphone + tag \texttt{v0.4} sur GitHub \\
\textbf{Rapport} & 1-2 pages : ce que j'ai fait + difficultés + lien repo \\
\bottomrule
\end{tabular}
\end{center}
\end{prereqbox}

\begin{objectifbox}[Objectifs du Lab]
À la fin de ce Lab, vous aurez :
\begin{enumerate}[nosep]
    \item Compris pourquoi votre app actuelle est \textbf{inutilisable sur téléphone}
    \item Ajouté la balise \texttt{<meta viewport>} pour le mobile
    \item Découvert et utilisé les \textbf{media queries} CSS
    \item Adopté l'approche \textbf{mobile-first}
    \item Adapté la \textbf{nav}, les \textbf{cartes}, et le \textbf{tableau} pour mobile
    \item Ajouté une \textbf{grille de drapeaux} responsive sur la page À propos
    \item Testé votre app dans le \textbf{mode mobile de DevTools}
    \item Posé le tag \texttt{v0.4} sur GitHub
\end{enumerate}
\end{objectifbox}

\tableofcontents

\newpage

% ===========================================================================
\section{Le problème : votre app n'est pas mobile}
% ===========================================================================

\subsection{Faites le test maintenant}

Ouvrez votre Jàng Afrig dans le navigateur. Maintenant :
\begin{enumerate}[nosep]
    \item Appuyez sur \texttt{F12} (ou \texttt{Cmd+Opt+I} sur Mac) pour ouvrir DevTools
    \item Cliquez sur l'icône \textbf{téléphone/tablette} en haut à gauche du panneau (ou \texttt{Ctrl+Shift+M})
    \item Choisissez un appareil dans la liste : \textbf{iPhone SE} ou \textbf{Pixel 5}
\end{enumerate}

Vous voyez l'horreur : la nav déborde, les cartes sont écrasées ou
minuscules, le tableau du palmarès dépasse l'écran à droite et oblige à
scroller horizontalement.

\subsection{Pourquoi est-ce si moche ?}

Deux raisons :

\begin{enumerate}[leftmargin=*,nosep]
    \item Vous n'avez pas la balise \texttt{<meta viewport>} dans vos pages, donc le téléphone affiche votre site comme s'il avait 980 pixels de large (et zoome dehors).
    \item Vous n'avez écrit \textbf{aucune règle CSS spécifique au mobile}. Tout est calculé pour un écran d'ordinateur (1100px de large dans votre \texttt{main}).
\end{enumerate}

\subsection{Ce qu'on va faire aujourd'hui}

\begin{itemize}[leftmargin=*,nosep]
    \item Ajouter le \texttt{<meta viewport>} (10 secondes)
    \item Adopter une approche \textbf{mobile-first} : on écrit d'abord pour téléphone, puis on adapte pour grand écran avec des \textbf{media queries}
    \item Adapter nav, cartes, tableau, et ajouter une grille de drapeaux
\end{itemize}

\begin{retenirbox}[Statistique mondiale]
En 2025, plus de \textbf{60\%} du trafic web mondial vient des téléphones.
Au Sénégal, ce chiffre dépasse \textbf{80\%}. Si votre site n'est pas
\emph{responsive}, vous abandonnez 8 utilisateurs sur 10. C'est
inacceptable.
\end{retenirbox}

\newpage

% ===========================================================================
\section{Le viewport meta tag (la ligne magique)}
% ===========================================================================

\subsection{Ajouter la ligne dans les 5 pages}

Dans le \texttt{<head>} de \textbf{chaque} page HTML, ajoutez cette ligne
juste après \texttt{<meta charset="UTF-8">} :

\begin{lstlisting}[language=HTML5]
<meta name="viewport" content="width=device-width, initial-scale=1">
\end{lstlisting}

\textbf{Enregistrez les 5 fichiers}. Rafraîchissez votre simulation
mobile dans DevTools : la page n'est plus dézoomée. Les éléments ont
maintenant leur taille \og{}réelle\fg{}. \cmark

\subsection{Ce que fait cette ligne}

\begin{itemize}[leftmargin=*,nosep]
    \item \texttt{width=device-width} : la largeur de la page égale celle de l'écran réel du téléphone (par exemple 375px sur iPhone SE)
    \item \texttt{initial-scale=1} : pas de zoom au chargement
\end{itemize}

\begin{infobox}[Avant/après cette ligne]
\textbf{Sans} \texttt{<meta viewport>} : le téléphone simule un écran de
980px et dézoome pour faire tenir la page. Résultat : texte minuscule.

\textbf{Avec} \texttt{<meta viewport>} : la page s'affiche à la vraie
taille de l'écran. Résultat : texte lisible mais\ldots la nav déborde
toujours. Il faut maintenant écrire des règles CSS pour mobile.
\end{infobox}

\begin{attentionbox}[À ne jamais oublier]
Mettez \texttt{<meta viewport>} sur \textbf{toutes} les pages HTML que
vous créerez dans votre vie. C'est la première ligne après
\texttt{<meta charset>}. C'est devenu un automatisme pour tous les
développeurs web depuis 2010.
\end{attentionbox}

\newpage

% ===========================================================================
\section{Mobile-first : la philosophie}
% ===========================================================================

\subsection{Deux approches possibles}

\begin{itemize}[leftmargin=*,nosep]
    \item \textbf{Desktop-first} : on écrit le CSS pour grand écran, puis on adapte pour mobile.
    \item \textbf{Mobile-first} : on écrit le CSS pour mobile, puis on adapte pour grand écran.
\end{itemize}

\subsection{Pourquoi mobile-first ?}

\begin{enumerate}[leftmargin=*,nosep]
    \item Le mobile est plus \textbf{contraint} (petit écran, peu d'espace). Si on règle d'abord ce cas difficile, ajouter de l'espace pour grand écran est facile.
    \item Le \textbf{layout par défaut} qu'on écrit (sans media query) est celui qui sera vu par les téléphones. C'est le cas le plus fréquent.
    \item Les \textbf{media queries de type \texttt{min-width}} sont plus lisibles que \texttt{max-width}.
\end{enumerate}

\subsection{Les 3 breakpoints qu'on va utiliser}

\begin{figure}[!ht]
\centering
\begin{tikzpicture}[
    every node/.style={font=\small\bfseries, anchor=north},
    bp/.style={rounded corners=4pt, draw, thick, minimum height=2.4cm, align=center}
]
    % Mobile
    \filldraw[fill=bpmobile!20, draw=bpmobile!70!black, thick, rounded corners=4pt]
        (0, 0) rectangle (3, 2.4);
    \node[anchor=north] at (1.5, 2.3) {\small Mobile};
    \node[anchor=center, font=\scriptsize] at (1.5, 1.6) {< 768px};
    \node[anchor=center, font=\scriptsize\itshape] at (1.5, 1.1) {(téléphone)};
    \node[anchor=center, font=\scriptsize] at (1.5, 0.4) {\textbf{défaut}};

    % Tablette
    \filldraw[fill=bptablet!30, draw=bptablet!70!black, thick, rounded corners=4pt]
        (3.5, 0) rectangle (8.5, 2.4);
    \node[anchor=north] at (6, 2.3) {\small Tablette};
    \node[anchor=center, font=\scriptsize] at (6, 1.6) {≥ 768px};
    \node[anchor=center, font=\scriptsize\itshape] at (6, 1.1) {(iPad, paysage mobile)};
    \node[anchor=center, font=\scriptsize\ttfamily] at (6, 0.4) {@media (min-width: 768px)};

    % Desktop
    \filldraw[fill=bpdesktop!20, draw=bpdesktop!70!black, thick, rounded corners=4pt]
        (9, 0) rectangle (15, 2.4);
    \node[anchor=north] at (12, 2.3) {\small Desktop};
    \node[anchor=center, font=\scriptsize] at (12, 1.6) {≥ 1024px};
    \node[anchor=center, font=\scriptsize\itshape] at (12, 1.1) {(ordinateur)};
    \node[anchor=center, font=\scriptsize\ttfamily] at (12, 0.4) {@media (min-width: 1024px)};
\end{tikzpicture}
\caption{Les 3 breakpoints utilisés dans ce Lab. Le CSS de base (sans media query) est pour mobile.}
\label{fig:breakpoints}
\end{figure}

\begin{retenirbox}[La règle mentale]
\begin{itemize}[nosep]
    \item Sans aucune media query $\to$ s'applique à \textbf{tous} les écrans (donc en mobile-first, c'est mon style mobile)
    \item \texttt{@media (min-width: 768px) \{\ldots\}} $\to$ s'ajoute pour les écrans \textbf{plus larges que 768px}
    \item \texttt{@media (min-width: 1024px) \{\ldots\}} $\to$ s'ajoute pour les écrans \textbf{plus larges que 1024px}
\end{itemize}
Plus l'écran est large, plus de règles s'empilent en se cumulant.
\end{retenirbox}

\newpage

% ===========================================================================
\section{Anatomie d'une media query}
% ===========================================================================

Une media query est une \textbf{condition} : si l'écran respecte une
certaine largeur, alors le CSS à l'intérieur s'applique.

\subsection{Syntaxe}

\begin{lstlisting}[language=CSS]
@media (min-width: 768px) {
    /* Toutes les règles ici ne s'appliquent que */
    /* si l'écran fait au moins 768 pixels de large */

    .carte-mode {
        flex: 1 1 250px;
    }

    main {
        padding: 0 40px;
    }
}
\end{lstlisting}

\subsection{Exemple complet : un titre qui change de taille}

\begin{lstlisting}[language=CSS]
/* Mobile (par défaut) */
header h1 {
    font-size: 1.8rem;
}

/* Tablette et plus */
@media (min-width: 768px) {
    header h1 {
        font-size: 2.2rem;
    }
}

/* Desktop et plus */
@media (min-width: 1024px) {
    header h1 {
        font-size: 2.8rem;
    }
}
\end{lstlisting}

\textbf{Comment lire ce code :}
\begin{itemize}[leftmargin=*,nosep]
    \item Sur téléphone (< 768px) : le titre est à 1.8rem
    \item Sur tablette (768-1023px) : le titre devient 2.2rem
    \item Sur desktop ($\geq$ 1024px) : le titre devient 2.8rem
\end{itemize}

\begin{infobox}[Pourquoi on n'écrit pas \texttt{max-width} ?]
On pourrait écrire \texttt{@media (max-width: 767px) \{\ldots\}} pour
viser uniquement le mobile. Mais en mobile-first, on préfère :
\begin{enumerate}[nosep]
    \item Écrire le mobile \emph{sans condition} (plus simple)
    \item \emph{Ajouter} pour les écrans plus grands avec \texttt{min-width}
\end{enumerate}
On a moins de doublons et le code est plus facile à lire.
\end{infobox}

\newpage

% ===========================================================================
\section{Étape 1 : Adapter le main pour mobile}
% ===========================================================================

Dans votre \texttt{style.css} actuel, vous avez probablement :

\begin{lstlisting}[language=CSS]
main {
    max-width: 1100px;
    margin: 30px auto;
    padding: 0 20px;
}
\end{lstlisting}

Le \texttt{max-width: 1100px} est OK pour mobile (la largeur sera celle
de l'écran), mais le \texttt{padding: 0 20px} peut être trop sur petit
écran. \textbf{Modifiez en} :

\begin{lstlisting}[language=CSS]
main {
    max-width: 1100px;
    margin: 20px auto;
    padding: 0 15px;
}

@media (min-width: 768px) {
    main {
        padding: 0 30px;
    }
}

@media (min-width: 1024px) {
    main {
        padding: 0 40px;
    }
}
\end{lstlisting}

\textbf{Enregistrez}. La marge intérieure du \texttt{main} grandit avec
l'écran. Sur mobile : 15px. Sur tablette : 30px. Sur desktop : 40px. \cmark

% ===========================================================================
\section{Étape 2 : Adapter la nav}
% ===========================================================================

Votre nav actuelle utilise \texttt{display: flex} avec
\texttt{justify-content: center}. Sur mobile, les 4 liens débordent.
\textbf{Solution simple :} les empiler en colonne sur petit écran, en
ligne sur grand écran.

\textbf{Remplacez} le bloc \texttt{nav} de votre CSS par :

\begin{lstlisting}[language=CSS]
/* ===== Navigation (mobile par défaut) ===== */
nav {
    background-color: var(--jaune);
    display: flex;
    flex-direction: column;
    text-align: center;
}

nav a {
    color: var(--noir);
    text-decoration: none;
    font-weight: 600;
    padding: 12px 15px;
    transition: background-color 0.2s;
    border-bottom: 1px solid rgba(0, 0, 0, 0.1);
}

nav a:hover {
    background-color: var(--vert);
    color: white;
}

/* Tablette et plus : on revient en ligne */
@media (min-width: 768px) {
    nav {
        flex-direction: row;
        justify-content: center;
        flex-wrap: wrap;
    }

    nav a {
        border-bottom: none;
        padding: 15px 25px;
    }
}
\end{lstlisting}

\textbf{Enregistrez}. Testez en redimensionnant la fenêtre :
\begin{itemize}[nosep]
    \item Petit écran : les liens sont \textbf{empilés verticalement}
    \item Grand écran : ils reviennent \textbf{en ligne horizontale}
\end{itemize}

\cmark

\begin{infobox}[Le hamburger menu, ce sera pour plus tard]
Sur les \emph{vrais} sites, on remplace souvent la nav mobile par une
icône \og{}hamburger\fg{} (☰) qui ouvre le menu au clic. Cela demande du
JavaScript, donc on s'en occupera au-delà de ce cours. Pour le L1, on
fait la version simple : les liens en colonne. C'est tout à fait
acceptable.
\end{infobox}

\newpage

% ===========================================================================
\section{Étape 3 : Adapter les cartes}
% ===========================================================================

Vos cartes \texttt{.carte-mode} ont \texttt{flex: 1 1 200px}, donc elles
se mettent déjà en grille flexible. Sur mobile, le \og{}1 1 200px\fg{}
peut donner 1 carte par ligne (parce que 200px est plus large que la
moitié de l'écran). Sur tablette, 2 cartes par ligne. Sur desktop, 3 ou
plus.

C'est déjà bien. Mais on peut affiner pour avoir un comportement plus
prévisible :

\textbf{Modifiez} \texttt{.carte-mode} dans votre CSS :

\begin{lstlisting}[language=CSS]
/* ===== Cartes de modes (mobile : 1 colonne) ===== */
.modes-grille {
    display: flex;
    flex-direction: column;
    gap: 12px;
}

.carte-mode {
    padding: 16px;
    border-radius: 8px;
    color: var(--noir);
    box-shadow: 0 2px 6px rgba(0, 0, 0, 0.1);
}

/* Tablette : 2 colonnes */
@media (min-width: 768px) {
    .modes-grille {
        flex-direction: row;
        flex-wrap: wrap;
    }
    .carte-mode {
        flex: 1 1 calc(50% - 12px);
    }
}

/* Desktop : 3 colonnes */
@media (min-width: 1024px) {
    .carte-mode {
        flex: 1 1 calc(33.333% - 12px);
    }
}
\end{lstlisting}

\textbf{Enregistrez}. Maintenant :
\begin{itemize}[nosep]
    \item Mobile : \textbf{1 carte par ligne}
    \item Tablette : \textbf{2 cartes par ligne}
    \item Desktop : \textbf{3 cartes par ligne}
\end{itemize}

\cmark

\begin{infobox}[\texttt{calc()}, c'est quoi ?]
\texttt{calc(50\% - 12px)} permet de \emph{calculer} une valeur en mêlant
unités. Ici : la moitié de la largeur du parent, moins 12 pixels (pour
compenser le \texttt{gap}). C'est très utile pour les layouts.
\end{infobox}

\newpage

% ===========================================================================
\section{Étape 4 : Une grille de drapeaux sur la page À propos}
% ===========================================================================

Sur \texttt{apropos.html}, ajoutez une nouvelle section qui affiche une
\textbf{galerie de drapeaux africains}.

\subsection{Récupérer les images}

Téléchargez 6 drapeaux (clic droit $\to$ \textbf{Enregistrer la cible
sous}) dans votre dossier \texttt{images/} :

\begin{itemize}[leftmargin=*,nosep]
    \item Sénégal : \url{https://flagcdn.com/w320/sn.png} → \texttt{sn.png}
    \item Mali : \url{https://flagcdn.com/w320/ml.png} → \texttt{ml.png}
    \item Maroc : \url{https://flagcdn.com/w320/ma.png} → \texttt{ma.png}
    \item Côte d'Ivoire : \url{https://flagcdn.com/w320/ci.png} → \texttt{ci.png}
    \item Égypte : \url{https://flagcdn.com/w320/eg.png} → \texttt{eg.png}
    \item Afrique du Sud : \url{https://flagcdn.com/w320/za.png} → \texttt{za.png}
\end{itemize}

\subsection{HTML : la grille}

\textbf{Ajoutez à \texttt{apropos.html}}, dans le \texttt{<main>} :

\begin{lstlisting}[language=HTML5]
<section>
    <h2>Quelques pays au programme</h2>
    <div class="grille-drapeaux">
        <figure>
            <img src="images/sn.png" alt="Drapeau du Sénégal">
            <figcaption>Sénégal</figcaption>
        </figure>
        <figure>
            <img src="images/ml.png" alt="Drapeau du Mali">
            <figcaption>Mali</figcaption>
        </figure>
        <figure>
            <img src="images/ma.png" alt="Drapeau du Maroc">
            <figcaption>Maroc</figcaption>
        </figure>
        <figure>
            <img src="images/ci.png" alt="Drapeau de la Côte d'Ivoire">
            <figcaption>Côte d'Ivoire</figcaption>
        </figure>
        <figure>
            <img src="images/eg.png" alt="Drapeau de l'Égypte">
            <figcaption>Égypte</figcaption>
        </figure>
        <figure>
            <img src="images/za.png" alt="Drapeau de l'Afrique du Sud">
            <figcaption>Afrique du Sud</figcaption>
        </figure>
    </div>
</section>
\end{lstlisting}

\subsection{CSS : la grille responsive}

\textbf{Ajoutez à \texttt{style.css}} :

\begin{lstlisting}[language=CSS]
/* ===== Grille de drapeaux (mobile : 2 par ligne) ===== */
.grille-drapeaux {
    display: flex;
    flex-wrap: wrap;
    gap: 15px;
    margin-top: 15px;
}

.grille-drapeaux figure {
    flex: 1 1 calc(50% - 15px);
    text-align: center;
    background-color: white;
    padding: 10px;
    border-radius: 6px;
    box-shadow: 0 2px 4px rgba(0, 0, 0, 0.08);
}

.grille-drapeaux img {
    width: 100%;
    height: auto;
    border-radius: 4px;
}

.grille-drapeaux figcaption {
    margin-top: 6px;
    font-size: 0.9rem;
    font-weight: 600;
}

/* Tablette : 3 par ligne */
@media (min-width: 768px) {
    .grille-drapeaux figure {
        flex: 1 1 calc(33.333% - 15px);
    }
}

/* Desktop : 6 par ligne */
@media (min-width: 1024px) {
    .grille-drapeaux figure {
        flex: 1 1 calc(16.666% - 15px);
    }
}
\end{lstlisting}

\textbf{Enregistrez}. Sur la page \texttt{apropos.html}, vous voyez les
6 drapeaux disposés :
\begin{itemize}[nosep]
    \item Mobile : 2 par ligne
    \item Tablette : 3 par ligne
    \item Desktop : 6 par ligne (alignés)
\end{itemize}

\cmark

\begin{infobox}[\texttt{<figure>} et \texttt{<figcaption>}]
Ces deux balises HTML5 sémantiques servent à associer une image à sa
légende. Plus propre que d'enchaîner \texttt{<img>} et \texttt{<p>}
anonymes : un lecteur d'écran annoncera \og{}figure : drapeau du Sénégal,
légende Sénégal\fg{}.
\end{infobox}

\newpage

% ===========================================================================
\section{Étape 5 : Adapter le tableau du palmarès}
% ===========================================================================

Le tableau a 5 colonnes (Rang, Pseudo, Mode, Score, Date). Sur mobile,
ça déborde et il faut scroller horizontalement. Plusieurs solutions
existent. La plus simple : permettre le \textbf{scroll horizontal
contrôlé} dans un conteneur.

\subsection{Modifier scores.html}

\textbf{Enveloppez} le tableau dans un \texttt{<div>} :

\begin{lstlisting}[language=HTML5]
<div class="tableau-conteneur">
    <table>
        <!-- ... lignes du tableau ... -->
    </table>
</div>
\end{lstlisting}

\subsection{Ajouter le CSS}

\begin{lstlisting}[language=CSS]
/* ===== Conteneur de tableau (scroll horizontal sur mobile) ===== */
.tableau-conteneur {
    overflow-x: auto;
    -webkit-overflow-scrolling: touch;
    margin-top: 15px;
}

table {
    min-width: 480px;
}

/* Police plus petite sur mobile */
@media (max-width: 767px) {
    th, td {
        font-size: 0.85rem;
        padding: 8px;
    }
}
\end{lstlisting}

\textbf{Enregistrez}. Sur mobile, le tableau a maintenant un
\textbf{ascenseur horizontal local} (juste pour le tableau, pas pour
toute la page). C'est élégant. \cmark

\begin{retenirbox}[Pourquoi cette approche ?]
Le tableau du palmarès a 5 colonnes \emph{essentielles}. On ne peut pas
en cacher sans perdre de l'info. Le scroll horizontal localisé est la
solution la plus simple et la plus courante. Plus tard, en L2, vous
verrez d'autres techniques (cartes par ligne, transformation en accordéon).
\end{retenirbox}

\newpage

% ===========================================================================
\section{Tester avec DevTools}
% ===========================================================================

\subsection{Le mode mobile de DevTools}

Dans Chrome ou Firefox :
\begin{enumerate}[nosep]
    \item Ouvrez DevTools (\texttt{F12})
    \item Cliquez sur l'icône \textbf{téléphone/tablette} en haut à gauche
    \item Sélectionnez un appareil dans la liste déroulante :
        \begin{itemize}[nosep]
            \item iPhone SE (375×667) --- petit téléphone
            \item iPhone 14 Pro (393×852) --- téléphone moderne
            \item iPad Mini (768×1024) --- tablette
            \item Responsive (taille libre, redimensionnez à la souris)
        \end{itemize}
\end{enumerate}

\subsection{Tester les 3 breakpoints}

\begin{enumerate}[nosep]
    \item Choisissez \textbf{iPhone SE} (375px) : tout doit s'afficher en colonne
    \item Choisissez \textbf{iPad Mini} (768px) : la nav passe en ligne, les cartes en 2 colonnes
    \item Passez en \textbf{Responsive} et étirez à 1100px : 3 cartes par ligne, 6 drapeaux par ligne
\end{enumerate}

\begin{infobox}[Tester sur un vrai téléphone (bonus)]
Si vous voulez voir Jàng Afrig sur \emph{votre} téléphone réel :
\begin{enumerate}[nosep]
    \item Connectez votre PC et votre téléphone au \textbf{même Wi-Fi}
    \item Dans VS Code, faites tourner Live Server
    \item Notez l'adresse IP de votre PC (\texttt{ipconfig} sur Windows, \texttt{ifconfig} sur Mac/Linux)
    \item Sur votre téléphone, ouvrez le navigateur et tapez : \texttt{http://192.168.X.X:5500/index.html}
\end{enumerate}
Ça marche, c'est très satisfaisant.
\end{infobox}

\subsection{Aperçu visuel attendu}

\begin{center}
\begin{tikzpicture}[
    every node/.style={anchor=north west, inner sep=4pt}
]
    % Mobile (375px)
    \draw[draw=gray, thick, rounded corners=8pt] (0, 0) rectangle (3.5, 6);
    \filldraw[fill=panvert, draw=panvert] (0.05, 5.5) rectangle (3.45, 5.95);
    \node[font=\tiny\bfseries, color=white, anchor=center] at (1.75, 5.725) {Jàng Afrig};

    \filldraw[fill=panjaune, draw=panjaune!70!black] (0.05, 4.4) rectangle (3.45, 5.5);
    \node[font=\tiny, anchor=center] at (1.75, 5.275) {Accueil};
    \node[font=\tiny, anchor=center] at (1.75, 5) {Jouer};
    \node[font=\tiny, anchor=center] at (1.75, 4.725) {Mes Scores};
    \node[font=\tiny, anchor=center] at (1.75, 4.5) {À propos};

    \filldraw[fill=panjaune!30, draw=panjaune!70!black, thick, rounded corners=2pt] (0.2, 3.5) rectangle (3.3, 4.2);
    \node[font=\tiny\bfseries, anchor=center] at (1.75, 3.85) {Capitales};

    \filldraw[fill=panrouge!30, draw=panrouge!70!black, thick, rounded corners=2pt] (0.2, 2.7) rectangle (3.3, 3.4);
    \node[font=\tiny\bfseries, color=white, anchor=center] at (1.75, 3.05) {Drapeaux};

    \filldraw[fill=panvert!30, draw=panvert!70!black, thick, rounded corners=2pt] (0.2, 1.9) rectangle (3.3, 2.6);
    \node[font=\tiny\bfseries, color=white, anchor=center] at (1.75, 2.25) {Langues};

    \node[anchor=center, font=\tiny\bfseries, color=bpmobile!70!black] at (1.75, -0.3) {Mobile (< 768px)};

    % Tablette (768px)
    \draw[draw=gray, thick, rounded corners=8pt] (4.5, 0) rectangle (10, 6);
    \filldraw[fill=panvert, draw=panvert] (4.55, 5.5) rectangle (9.95, 5.95);
    \node[font=\tiny\bfseries, color=white, anchor=center] at (7.25, 5.725) {Jàng Afrig};

    \filldraw[fill=panjaune, draw=panjaune!70!black] (4.55, 5) rectangle (9.95, 5.5);
    \node[font=\tiny, anchor=center] at (5.4, 5.25) {Accueil};
    \node[font=\tiny, anchor=center] at (6.6, 5.25) {Jouer};
    \node[font=\tiny, anchor=center] at (7.85, 5.25) {Mes Scores};
    \node[font=\tiny, anchor=center] at (9.1, 5.25) {À propos};

    \filldraw[fill=panjaune!30, draw=panjaune!70!black, thick, rounded corners=2pt] (4.7, 4) rectangle (7.2, 4.8);
    \node[font=\tiny\bfseries, anchor=center] at (5.95, 4.4) {Capitales};
    \filldraw[fill=panrouge!30, draw=panrouge!70!black, thick, rounded corners=2pt] (7.3, 4) rectangle (9.8, 4.8);
    \node[font=\tiny\bfseries, color=white, anchor=center] at (8.55, 4.4) {Drapeaux};

    \filldraw[fill=panvert!30, draw=panvert!70!black, thick, rounded corners=2pt] (4.7, 3.1) rectangle (7.2, 3.9);
    \node[font=\tiny\bfseries, color=white, anchor=center] at (5.95, 3.5) {Langues};
    \filldraw[fill=panjaune!30, draw=panjaune!70!black, thick, rounded corners=2pt] (7.3, 3.1) rectangle (9.8, 3.9);
    \node[font=\tiny\bfseries, anchor=center] at (8.55, 3.5) {Voisins};

    \node[anchor=center, font=\tiny\bfseries, color=bptablet!70!black] at (7.25, -0.3) {Tablette ($\geq$ 768px)};

    % Desktop (1100px)
    \draw[draw=gray, thick, rounded corners=8pt] (11, 0) rectangle (16.5, 6);
    \filldraw[fill=panvert, draw=panvert] (11.05, 5.5) rectangle (16.45, 5.95);
    \node[font=\tiny\bfseries, color=white, anchor=center] at (13.75, 5.725) {Jàng Afrig};
    \filldraw[fill=panjaune, draw=panjaune!70!black] (11.05, 5) rectangle (16.45, 5.5);
    \node[font=\tiny, anchor=center] at (12, 5.25) {Accueil};
    \node[font=\tiny, anchor=center] at (13.2, 5.25) {Jouer};
    \node[font=\tiny, anchor=center] at (14.4, 5.25) {Scores};
    \node[font=\tiny, anchor=center] at (15.6, 5.25) {À propos};

    \filldraw[fill=panjaune!30, draw=panjaune!70!black, thick, rounded corners=2pt] (11.2, 4) rectangle (12.85, 4.8);
    \node[font=\tiny\bfseries, anchor=center] at (12, 4.4) {Capitales};
    \filldraw[fill=panrouge!30, draw=panrouge!70!black, thick, rounded corners=2pt] (12.95, 4) rectangle (14.6, 4.8);
    \node[font=\tiny\bfseries, color=white, anchor=center] at (13.78, 4.4) {Drapeaux};
    \filldraw[fill=panvert!30, draw=panvert!70!black, thick, rounded corners=2pt] (14.7, 4) rectangle (16.35, 4.8);
    \node[font=\tiny\bfseries, color=white, anchor=center] at (15.5, 4.4) {Langues};

    \node[anchor=center, font=\tiny\bfseries, color=bpdesktop!70!black] at (13.75, -0.3) {Desktop ($\geq$ 1024px)};
\end{tikzpicture}
\end{center}

\textbf{La même page} affichée différemment selon la largeur. C'est ça,
le responsive.

\newpage

% ===========================================================================
\section{Exercices d'application}
% ===========================================================================

\begin{objectifbox}[Trois mini-exercices]

\textbf{Exercice 1 --- Footer responsive}

Faites en sorte que sur mobile, le footer affiche le copyright sur une
ligne et le lien \og{}En savoir plus\fg{} sur la ligne suivante (au lieu
de tout sur une ligne qui déborde). Indice : \texttt{display: flex;
flex-direction: column; gap: 8px;}.

\medskip

\textbf{Exercice 2 --- Image héro responsive}

Sur la page d'accueil (\texttt{index.html}), ajoutez juste après le
\texttt{<header>} une grande image \og{}héro\fg{} (par exemple un drapeau
de l'Union Africaine).

\begin{itemize}[nosep]
    \item Largeur 100\% de l'écran
    \item Hauteur max 200px sur mobile, 300px sur desktop
    \item Centrée, avec \texttt{object-fit: cover}
\end{itemize}

\medskip

\textbf{Exercice 3 --- Formulaire responsive}

Sur \texttt{quiz.html}, faites en sorte que le formulaire occupe au
maximum :
\begin{itemize}[nosep]
    \item \textbf{100\%} de la largeur sur mobile
    \item \textbf{500px} centrés sur tablette et desktop
\end{itemize}

Indice : \texttt{form \{ max-width: 500px; margin: 0 auto; \}}.

\medskip

\textbf{Bonus} --- Affichage \og{}sticky\fg{}

Faites en sorte que la nav reste \textbf{collée en haut} de l'écran quand
on scrolle. Une seule ligne de CSS suffit : \texttt{position: sticky;
top: 0;} sur \texttt{nav}.
\end{objectifbox}

\newpage

% ===========================================================================
\section{Aide-mémoire}
% ===========================================================================

\subsection{Le b.a.-ba du responsive}

\begin{center}
\begin{tabular}{@{}lp{8cm}@{}}
\toprule
\textbf{Élément}                                & \textbf{Rôle} \\
\midrule
\texttt{<meta name="viewport" content="...">}   & La ligne magique pour le mobile \\
\addlinespace
\texttt{@media (min-width: 768px)}              & Règles pour tablette et plus \\
\texttt{@media (min-width: 1024px)}             & Règles pour desktop et plus \\
\addlinespace
\texttt{flex-direction: column / row}           & Empiler / aligner \\
\texttt{flex-wrap: wrap}                        & Retour à la ligne \\
\texttt{flex: 1 1 calc(50\% - 15px)}            & 2 colonnes flexibles \\
\addlinespace
\texttt{overflow-x: auto}                       & Scroll horizontal local \\
\texttt{position: sticky; top: 0;}              & Élément collé en haut au scroll \\
\bottomrule
\end{tabular}
\end{center}

\subsection{Les 3 breakpoints}

\begin{center}
\begin{tabular}{@{}lll@{}}
\toprule
\textbf{Plage}        & \textbf{Cible}        & \textbf{Media query} \\
\midrule
< 768px               & Mobile (par défaut) & --- (CSS de base) \\
768 -- 1023px         & Tablette / paysage   & \texttt{@media (min-width: 768px)} \\
$\geq$ 1024px         & Desktop              & \texttt{@media (min-width: 1024px)} \\
\bottomrule
\end{tabular}
\end{center}

\subsection{Méthodologie mobile-first}

\begin{enumerate}[leftmargin=*,nosep]
    \item Écrire d'abord le CSS pour mobile (\textbf{sans} media query)
    \item Tester sur DevTools mode mobile (iPhone SE 375px)
    \item Ajouter \texttt{@media (min-width: 768px)} pour ce qui doit changer en tablette
    \item Tester en mode tablette (iPad Mini 768px)
    \item Ajouter \texttt{@media (min-width: 1024px)} pour le desktop
    \item Tester en mode responsive à 1100px
\end{enumerate}

\subsection{Tester sur vrai téléphone}

\begin{enumerate}[nosep]
    \item PC et téléphone sur le même Wi-Fi
    \item \texttt{ipconfig} (Windows) / \texttt{ifconfig} (Mac/Linux) pour récupérer l'IP du PC
    \item Sur téléphone : \texttt{http://IP-DU-PC:5500/index.html}
\end{enumerate}

\newpage

% ===========================================================================
\section{Routine Git}
% ===========================================================================

\begin{lstlisting}
git status
git add .
git commit -m "v0.4 : responsive mobile + grille drapeaux"
git push
git tag v0.4
git push --tags
\end{lstlisting}

Vérifiez sur \url{https://github.com/votre-username/jang-afrig} :
\begin{itemize}[nosep]
    \item Le \texttt{<meta viewport>} est dans les 5 pages
    \item \texttt{style.css} contient les media queries
    \item Les 6 drapeaux sont dans \texttt{images/}
    \item Le tag \texttt{v0.4} est visible
\end{itemize}

% ===========================================================================
\section{Rapport à produire}
% ===========================================================================

\begin{attentionbox}[Rapport du Lab 4 --- obligatoire]
\textbf{Format :} 1-2 pages, déposé dans \texttt{rapports/rapport-lab-4-NOM-Prenom.pdf}.

\medskip

\textbf{Contenu (4 sections) :}
\begin{enumerate}[nosep]
    \item \textbf{Ce que j'ai fait} : viewport ajouté, media queries écrites, nav/cartes/tableau adaptés, grille de drapeaux, exercices d'application.
    \item \textbf{Difficultés rencontrées} : layout qui casse, breakpoints qui ne marchent pas, scroll horizontal indésirable.
    \item \textbf{Comment j'ai résolu} : DevTools mode mobile, MDN, aide d'un camarade.
    \item \textbf{Lien vers mon repo avec le tag v0.4} : \\
          ex. \url{https://github.com/fatou-diop/jang-afrig/releases/tag/v0.4}
\end{enumerate}

\medskip

\textbf{Bonus dans le rapport :} \textbf{2 captures d'écran} de la même
page : une en mode mobile (375px), une en mode desktop (1100px). C'est
votre preuve fière que vous savez faire du responsive.
\end{attentionbox}

\newpage

% ===========================================================================
\section{Résumé et transition}
% ===========================================================================

\subsection{Ce que vous avez accompli}

\begin{itemize}[leftmargin=*,nosep]
    \item Ajouté \texttt{<meta viewport>} aux 5 pages \cmark
    \item Adopté l'approche \textbf{mobile-first} \cmark
    \item Maîtrisé les \textbf{media queries} avec 3 breakpoints \cmark
    \item Adapté nav, cartes, formulaire, tableau \cmark
    \item Construit une grille de drapeaux responsive \cmark
    \item Testé dans le mode mobile de DevTools \cmark
    \item Posé le tag \texttt{v0.4} \cmark
\end{itemize}

\subsection{Le bloc CSS est terminé}

Vous savez maintenant tout le CSS dont vous avez besoin pour ce semestre :
sélecteurs, cascade, box model, Flexbox, variables, media queries,
responsive. C'est solide.

\subsection{Ce qu'on fera au CM3 et au Lab 5}

\textbf{CM3 (S7) : JavaScript débarque.} Le langage qui anime les pages.
On verra variables, fonctions, conditions, boucles. Plus de pages
statiques --- on va commencer à \emph{réagir} aux clics de l'utilisateur.

\textbf{Lab 5 (S8) :} JS pratique. Vous découvrirez la console DevTools,
écrirez vos premières lignes de JS, et ferez bouger un élément de Jàng
Afrig avec un clic.

\begin{projetbox}[À ce stade du semestre]
Votre application a \textbf{une vraie interface professionnelle} qui
marche \textbf{sur n'importe quel appareil}. Plus de la moitié du chemin
est faite.

Manque maintenant le \textbf{cerveau} : la logique JavaScript qui rend
l'app \textit{interactive}. C'est ce qu'on attaque dans la prochaine
séance.
\end{projetbox}

\vspace{0.5cm}

\begin{center}
\fbox{
\begin{minipage}{0.85\textwidth}
\centering
\vspace{0.5cm}
\textbf{\Large Lab 4 --- Terminé \cmark}\\[0.3cm]
\normalsize
Votre app fonctionne sur tous les écrans, du téléphone à l'ordinateur.\\
Place à JavaScript : on va lui donner vie.\\[0.2cm]
\textit{Rendez-vous au CM3 !}\\[0.3cm]
\vspace{0.3cm}
\end{minipage}
}
\end{center}

\end{document}
